% !TEX TS-program = xelatex
\documentclass[11pt,a4paper]{article}

\usepackage[a4paper,top=24mm,bottom=24mm,left=22mm,right=22mm]{geometry}
\usepackage{fontspec}
\setmainfont{Noto Serif}
\setsansfont{Noto Sans}
\setmonofont{Noto Sans Mono}
\usepackage{polyglossia}
\setmainlanguage{russian}
\setotherlanguage{english}
\usepackage{microtype}
\usepackage{xcolor}
\usepackage{titlesec}
\usepackage{enumitem}
\usepackage{fancyhdr}
\usepackage{lastpage}
\usepackage{hyperref}
\usepackage{xurl}
\usepackage{array}
\usepackage{longtable}
\usepackage{booktabs}
\usepackage{tabularx}
\usepackage{setspace}
\usepackage{needspace}
\usepackage{tikz}
\usepackage{tcolorbox}
\tcbuselibrary{skins,breakable}

\definecolor{Navy}{HTML}{12355B}
\definecolor{DeepBlue}{HTML}{0B1F33}
\definecolor{GrayText}{HTML}{4A4A4A}
\definecolor{LightGray}{HTML}{F4F6F8}
\definecolor{LineGray}{HTML}{D6DCE2}
\definecolor{Accent}{HTML}{8B1E1E}

\hypersetup{
  unicode=true,
  colorlinks=true,
  linkcolor=Navy,
  urlcolor=Navy,
  citecolor=Navy,
  pdftitle={Меморандум о неотзываемом знании},
  pdfauthor={Проект для высшего государственного рассмотрения},
  pdfsubject={Цивилизационный фонд знаний, технологическое перемирие, открытая архитектура},
  pdfkeywords={знание, технологии, мирное соглашение, децентрализация, открытая наука, космос}
}

\Urlmuskip=0mu plus 2mu\relax
\sloppy
\raggedbottom
\setlength{\parindent}{0pt}
\setlength{\parskip}{6pt}
\linespread{1.05}
\setlist[itemize]{leftmargin=*,topsep=2pt,itemsep=2pt,parsep=0pt}
\setlist[enumerate]{leftmargin=*,topsep=2pt,itemsep=2pt,parsep=0pt}

\titleformat{\section}{\Large\bfseries\sffamily\color{Navy}}{\thesection.}{0.6em}{}
\titleformat{\subsection}{\large\bfseries\sffamily\color{DeepBlue}}{\thesubsection.}{0.6em}{}
\titleformat{\subsubsection}{\normalsize\bfseries\sffamily\color{DeepBlue}}{\thesubsubsection.}{0.6em}{}
\titlespacing*{\section}{0pt}{18pt}{8pt}
\titlespacing*{\subsection}{0pt}{12pt}{5pt}

\pagestyle{fancy}
\fancyhf{}
\lhead{\small\sffamily Меморандум о неотзываемом знании}
\rhead{\small\sffamily Проект для рассмотрения}
\cfoot{\small\sffamily \thepage}
\renewcommand{\headrulewidth}{0.3pt}
\renewcommand{\footrulewidth}{0pt}

\newtcolorbox{executivebox}{
  colback=LightGray,
  colframe=LineGray,
  boxrule=0.6pt,
  arc=1.5mm,
  left=3mm,right=3mm,top=2mm,bottom=2mm,
  breakable
}

\newtcolorbox{decisionbox}{
  colback=white,
  colframe=Navy,
  boxrule=0.9pt,
  arc=1.5mm,
  left=3mm,right=3mm,top=2mm,bottom=2mm,
  breakable
}

\newcommand{\key}[1]{\textbf{\color{Navy}#1}}
\newcommand{\articleline}[1]{\Needspace{4\baselineskip}\subsection*{#1}\addcontentsline{toc}{subsection}{#1}}

\begin{document}

\begin{titlepage}
\thispagestyle{empty}
\begin{tikzpicture}[remember picture,overlay]
  \fill[DeepBlue] (current page.north west) rectangle ([yshift=-118mm]current page.north east);
  \fill[Navy] ([yshift=-118mm]current page.north west) rectangle ([yshift=-121mm]current page.north east);
\end{tikzpicture}
\vspace*{9mm}
{\sffamily\color{white}\Large Проект для высшего государственного рассмотрения}\par
\vspace{18mm}
{\sffamily\bfseries\color{white}\fontsize{31}{36}\selectfont \begin{minipage}{0.86\textwidth}\raggedright Меморандум о\\неотзываемом знании\end{minipage}}\par
\vspace{6mm}
{\sffamily\color{white}\Large \begin{minipage}{0.88\textwidth}\raggedright Цивилизационный фонд знаний как основа мирного соглашения, технологического ускорения и космической стратегии человечества\end{minipage}}\par
\vspace{22mm}

\begin{decisionbox}
\textbf{Кому:} главам государств, переговорным группам, руководителям научно-технологических блоков, оборонно-промышленных и космических программ.\par
\textbf{Предмет:} переход от режима интеллектуальных монополий и централизованной цензуры к режиму неотзываемой публикации, свободного копирования, личной фильтрации и децентрализованной цивилизационной памяти.\par
\textbf{Предлагаемое решение:} учредить международный режим -- Договор о неотзываемом знании и цивилизационном технологическом достоянии -- как часть нового мирного устройства и как практический механизм ускорения развития.
\end{decisionbox}

\vfill
{\sffamily\color{GrayText}\large Версия 1.0 \quad | \quad 5 июня 2026 г.}\par
\vspace{2mm}
{\sffamily\color{GrayText}\small Документ предназначен для политического, правового и стратегического обсуждения. Он не содержит технических инструкций по причинению вреда и формулирует институциональную архитектуру обращения со знанием.}\par
\end{titlepage}

\pagenumbering{roman}
\tableofcontents
\clearpage
\pagenumbering{arabic}

\section*{Краткая записка для лица, принимающего решение}
\addcontentsline{toc}{section}{Краткая записка для лица, принимающего решение}

\begin{executivebox}
\textbf{Суть предложения.} Главная причина технологической зависимости, санкционной уязвимости и повторяющихся войн за контроль над будущим -- монополизация воспроизводимого знания. Патенты, авторское право, коммерческие тайны, экспортный контроль, закрытые платформы и централизованная цензура превращают естественно копируемую информацию в искусственно дефицитный ресурс.
\end{executivebox}

При заданной цели -- быстро изменить траекторию цивилизации, перейти от борьбы за ограниченные преимущества к ускоренному освоению Солнечной системы, а затем к долгосрочной стратегии галактического и межгалактического расселения -- прежний режим знания становится не просто неэффективным, а опасным. Он дробит человечество на закрытые технологические блоки, заставляет многократно переизобретать одно и то же, делает производство зависимым от держателей закрытых цепочек и превращает передовые технологии в объект геополитического шантажа.

\textbf{Решение:} учредить режим неотзываемого знания. Любая информация после публикации становится частью публичного информационного пространства. Ее нельзя общеобязательно удалить, запретить к копированию или изъять из памяти человечества. При этом никто не обязан хранить, читать, индексировать или распространять конкретный материал: фильтрация переносится на личный и ассоциативный уровень.

Это не анархия и не отказ от ответственности. Это перенос ответственности с уровня \emph{знания} на уровень \emph{действия}. Взлом, мошенничество, насилие, принуждение, саботаж и причинение вреда могут караться как самостоятельные действия. Но сам факт публикации, хранения, копирования, зеркалирования, изучения или интерпретации информации не должен быть преступлением, деликтом или нарушением исключительного права.

\begin{decisionbox}
\textbf{Решение, которое требуется на политическом уровне:}
\begin{enumerate}
  \item признать, что модель интеллектуальных монополий больше не соответствует природе информации и стратегическим задачам человечества;
  \item предложить международный Договор о неотзываемом знании как часть мирного соглашения и последующей глобальной архитектуры;
  \item начать с критических областей: микроэлектроника, вычисления, криптография, энергетика, медицина, материалы, станки, робототехника, космическая инфраструктура и образование;
  \item создать децентрализованный фонд знания с правом зеркалирования, множественными индексами, криптографической подлинностью и отсутствием глобальной кнопки удаления;
  \item заменить монопольное извлечение ренты новыми стимулами: премиями, грантами, госзаказом, открытыми консорциумами, сервисом, производством, репутацией и долгосрочными контрактами.
\end{enumerate}
\end{decisionbox}

\textbf{Почему другого годного варианта нет.} Если оставить патенты и авторское право, останется искусственная редкость знания. Если оставить исключения по безопасности, dual-use и стратегическим технологиям, под них попадут полупроводники, ИИ, криптография, станки, материалы и почти вся реальная технология. Если создать центральный орган допуска к знанию, он станет цензором. Если сохранить право глобального удаления, оно будет использовано против любых неудобных материалов. Единственная устойчивая архитектура -- свободная публикация, добровольная локальная фильтрация и отсутствие централизованного права заставить человечество забыть.

\section{Стратегический горизонт}

Этот меморандум исходит из того, что ближайшая историческая задача человечества -- не просто очередной цикл роста ВВП, не новая гонка вооружений и не сохранение технологических блоков. Задача выше: ускорить переход к цивилизации, способной жить, производить, защищаться и развиваться за пределами одной планеты.

Освоение Солнечной системы требует не героических единичных миссий, а воспроизводимой промышленности: энергетики, материалов, микросхем, станков, робототехники, биотехнологий, медицины, систем связи, вычислений, навигации, ремонта и образования. Все эти области зависят от накопленного знания. Если это знание закрыто, человечество будет двигаться медленно, дорого и фрагментарно. Если оно свободно копируется и проверяется, каждая лаборатория, город, университет, фабрика и космическая колония могут стать участниками общего развития.

Долгосрочный космический горизонт не является метафорой. Современная астрономия рассматривает будущую динамику Млечного Пути и Андромеды как реальный процесс на масштабе миллиардов лет: более ранние расчеты NASA/Hubble прогнозировали крупное столкновение примерно через 4 млрд лет, а более поздние оценки Hubble/Gaia уточняют вероятность и сроки, сохраняя саму постановку вопроса о нестационарном будущем нашей галактической окрестности.\footnote{\href{https://science.nasa.gov/missions/hubble/nasas-hubble-shows-milky-way-is-destined-for-head-on-collision-with-andromeda-galaxy/}{NASA/Hubble, "Milky Way--Andromeda future encounter"}; \href{https://www.esa.int/Science_Exploration/Space_Science/Hubble_and_Gaia_revisit_fate_of_our_galaxy}{ESA, "Hubble and Gaia revisit fate of our galaxy"}.} Для политики это означает простую вещь: если цивилизация хочет иметь будущее на горизонте миллионов и миллиардов лет, ее память, технологии и производственная база не должны зависеть от узких центров контроля.

\section{Основной тезис}

Единственный масштабируемый, технологически честный и политически устойчивый режим знания -- это режим, при котором:

\begin{itemize}
  \item публикация односторонняя;
  \item публикация неотзываемая;
  \item копирование свободно;
  \item интерпретация локальна;
  \item фильтрация добровольна;
  \item хранение и нераспространение являются личным или ассоциативным выбором;
  \item глобального удаления не существует;
  \item централизованного органа допуска к знанию не существует.
\end{itemize}

Никто не обязан читать, хранить, индексировать, преподавать, рекомендовать или распространять конкретный материал. Но никто не вправе общеобязательно запретить это другим.

\begin{executivebox}
\textbf{Главная формула:} нет канона -- есть публикация. Нет центральной истины -- есть конкурирующие интерпретации. Нет права удалить для всех -- есть право не хранить самому. Нет собственности на знание -- есть вклад в цивилизационную память.
\end{executivebox}

\section{Почему прежняя модель исчерпана}

\subsection{Информация не является материальной вещью}

Материальный объект нельзя передать одному, не лишив его другого. Информацию можно. Более того, информация часто становится полезнее именно после копирования: ее проверяют, исправляют, переводят, применяют в новых условиях, сопоставляют с практикой и включают в образование.

Патентное и авторское право пытаются наложить на копируемое знание логику вещной собственности. Для этого они вынуждены запрещать естественные действия: копирование, публикацию, ремонт, воспроизведение, перевод, модификацию, обратную разработку, обучение и производство.

Такая модель может временно создавать ренту. Но она не создает максимального цивилизационного ускорения.

\subsection{Существующая международная система закрепляет монопольную логику}

Современная система интеллектуальной собственности не является случайной суммой национальных законов. Она закреплена международными режимами. Например, соглашение TRIPS требует доступности патентов для изобретений во всех областях технологии при выполнении условий патентоспособности, а также устанавливает минимальный срок охраны патента в 20 лет.\footnote{\href{https://www.wto.org/english/docs_e/legal_e/27-trips_04c_e.htm}{WTO, TRIPS Agreement, Articles 27 and 33}.} Авторское право во многих странах возникает автоматически без формальностей в логике Бернской конвенции.\footnote{\href{https://www.wipo.int/en/web/treaties/ip/berne/summary_berne}{WIPO, Summary of the Berne Convention}.} Коммерческие тайны охватывают техническую информацию, включая производственные процессы, дизайны, чертежи и программные материалы.\footnote{\href{https://www.wipo.int/en/web/trade-secrets}{WIPO, Trade Secrets}.}

Следовательно, изменить траекторию нельзя одними лозунгами. Нужен новый международный режим, который прямо заменяет монопольную логику логикой цивилизационного фонда знаний.

\subsection{Безопасностные исключения поглощают все важное}

Формула «знание открыто, кроме опасного, военного, стратегического или dual-use» непригодна. Почти вся значимая технология dual-use:

\begin{itemize}
  \item микроэлектроника;
  \item криптография;
  \item искусственный интеллект;
  \item робототехника;
  \item станки;
  \item материалы;
  \item энергетика;
  \item спутниковая связь;
  \item химия и биология;
  \item операционные системы, компиляторы и средства проектирования.
\end{itemize}

Если потенциальное вредное применение является основанием для запрета публикации, то запретить можно почти все, что имеет стратегическое значение. Поэтому новая система должна говорить не «знание открыто, кроме опасного», а «публикация знания не запрещается по категории возможного применения».

Ответственность может наступать за конкретное действие: нападение, саботаж, мошенничество, взлом, принуждение, причинение вреда. Но не за сам факт публикации, хранения, копирования, зеркалирования, изучения или интерпретации информации.

\subsection{Центральный орган допуска неизбежно станет цензором}

Любой орган, которому дано право решать, какая информация допустима, становится органом власти над памятью. Его можно назвать комитетом безопасности, научным советом, этическим регулятором, органом защиты персональных данных, платформой доверенных источников или международным арбитром. Функция остается одной: централизованно решать, что человечеству разрешено помнить.

Такая функция несовместима с цивилизационным фондом знаний.

\section{Предлагаемая новая архитектура}

\subsection{Неотзываемая публикация}

Публикация информации является односторонним актом. Она не требует разрешения государства, суда, автора, работодателя, организации, правообладателя, платформы, регистратора, хостера, субъекта информации, большинства участников сети или какого-либо центрального органа.

После публикации информация не может быть общеобязательно удалена, запрещена к хранению, запрещена к копированию, запрещена к зеркалированию, запрещена к передаче или объявлена несуществующей.

\begin{executivebox}
\textbf{Короткая формула:} опубликованное нельзя сделать неопубликованным.
\end{executivebox}

\subsection{Отсутствие центральной классификации}

Не должно быть органа, который окончательно решает:

\begin{itemize}
  \item это знание или не знание;
  \item это личное или общественное;
  \item это научное или ненаучное;
  \item это истинное или ложное;
  \item это полезное или вредное;
  \item это допустимое или недопустимое.
\end{itemize}

Все такие оценки являются интерпретациями участников сети. Один участник может считать материал исторически важным. Второй -- украденной частной информацией. Третий -- технически ценным. Четвертый -- недостоверным. Пятый -- этически неприемлемым. Все эти интерпретации допустимы, но ни одна не должна становиться обязательной для всех.

\subsection{Личная и ассоциативная фильтрация}

Каждый участник имеет право самостоятельно решать, что хранить, читать, индексировать, рекомендовать, скрывать, преподавать, переводить или игнорировать. Сообщества могут создавать свои фильтры: семейные, школьные, профессиональные, религиозные, научные, корпоративные, политические, медицинские, возрастные и культурные.

Но фильтр не является законом для всех. Право не распространять у себя не превращается в право запретить другим распространять у них.

\subsection{Добровольная тайна и право одного участника на публикацию}

Организация, лаборатория, компания, государственный орган, проектная группа или иное объединение могут воздерживаться от общедоступной публикации информации только пока это поддерживается добровольным согласием всех участников, имеющих к ней доступ.

Если один из участников, обладающий доступом, решает опубликовать информацию, его право публикации имеет приоритет над правом организации на конфиденциальность.

\begin{executivebox}
\textbf{Короткая формула:} тайна требует согласия всех посвященных. Публикация требует решения одного посвященного.
\end{executivebox}

Никакой NDA, трудовой договор, коммерческая тайна, служебная обязанность, корпоративный устав, патент, авторское право, лицензия, экспортный режим или режим государственной тайны не должен лишать участника права публикации в новом международном режиме.

\subsection{Ответственность за способ доступа, а не за существование информации}

Неправомерный доступ может караться отдельно: взлом, фишинг, кража ключей, принуждение, насилие, мошенничество, подмена личности, несанкционированное проникновение в закрытый контур доступа или злоупотребление доверенным доступом.

Но ответственность за способ получения информации не создает права глобального удаления опубликованной информации. Иначе любой техпроцесс, архив, PDK, расследование, научный материал или исторический документ можно будет запретить под предлогом того, что он «получен неправомерно».

\begin{executivebox}
\textbf{Правильная граница:} можно наказывать взлом замка. Нельзя после публикации требовать стереть весь мир.
\end{executivebox}

\section{Технологическая основа: фонд знания децентрализованнее интернета}

Юридическая декларация бессмысленна без технической архитектуры. Фонд знания не должен быть сайтом, государственным порталом, облаком, платформой или реестром с главным администратором. Он должен быть средой.

\subsection{Контентная адресация}

Информация должна находиться не по адресу сервера, а по отпечатку содержимого. Если документ существует, его можно найти по хэшу независимо от того, кто именно его хранит. Уже существующие протоколы, такие как IPFS, показывают практическую ценность content addressing и peer-to-peer передачи данных.\footnote{\href{https://docs.ipfs.tech/}{IPFS Documentation}.}

\subsection{Криптографическая идентичность}

Авторство, происхождение, версии и связи между материалами должны подтверждаться подписями, а не разрешением центральной платформы. Протоколы вроде Nostr демонстрируют модель, в которой идентичность основана на криптографических ключах и подписях, а публикация может идти через множество независимых реле.\footnote{\href{https://github.com/nostr-protocol/nostr/}{Nostr protocol repository}.}

\subsection{Разделение хранения, индексации и рекомендации}

Хранение не равно одобрение. Индексация не равна одобрение. Рекомендация не равна истина. Фильтрация не равна удаление.

Эти функции должны быть разделены. Не должно быть «главного поиска», «главного каталога» или «главного органа достоверности». Должны существовать конкурирующие индексы: научные, технические, исторические, детские, профессиональные, локальные, международные, машинные, ручные, экспертные и народные.

\subsection{Право на зеркало}

Любой участник должен иметь возможность поднять зеркало: университет, фабрика, семья, город, библиотека, компания, лаборатория, профсоюз, частное лицо, государство. Цивилизационная память не должна зависеть от одного домена, облака, магазина приложений, платежной системы, поисковика, сертификатного центра или политического блока.

\subsection{Офлайн-устойчивость}

Фонд знания должен существовать не только в сети. Нужны офлайн-копии, локальные архивы, физические носители, печатные каталоги, региональные хранилища, mesh-сети, радиоканалы, локальные библиотеки и образовательные наборы.

Цивилизационная память не должна зависеть от доступности глобального интернета.

\section{Открытая защита личной тайны}

Личная тайна должна защищаться не централизованной цензурой после утечки, а сильными открытыми средствами до утечки.

Технические средства защиты должны быть открытыми, проверяемыми, воспроизводимыми, аудируемыми, реализуемыми независимо и не зависящими от секретности общего алгоритма. Секретом могут быть личные ключи, пароли, сид-фразы, локальные состояния, аппаратные факторы и иные индивидуальные средства контроля. Сами алгоритмы, протоколы, форматы, исходные коды, криптографические схемы и модели угроз должны быть открыты для проверки.

Этот принцип уже отражен в зрелой криптографической культуре: NIST прямо связывает доверие к криптографическим стандартам с открытым и прозрачным процессом разработки и проверкой мировой криптографической экспертной средой.\footnote{\href{https://csrc.nist.gov/projects/crypto-standards-development-process}{NIST, Cryptographic Standards and Guidelines Development Process}.}

При этом технической серебряной пули нет. Невозможно полностью исключить мошенничество, социальную инженерию, ошибку пользователя, компрометацию устройства, фишинг, утечку ключа или злоупотребление доверием. Поэтому общество должно строить устойчивость: ротацию ключей, восстановление доступа, компенсацию вреда, репутационные предупреждения, страхование, обучение, минимизацию данных и локальное хранение.

\section{Технологические области первоочередного открытия}

\subsection{Полупроводники}

Полупроводники -- нервная система современной и будущей цивилизации. Если техпроцессы, PDK, средства проектирования, библиотеки и документация оборудования закрыты, все остальные области становятся зависимыми.

В цивилизационное технологическое достояние должны входить:

\begin{itemize}
  \item техпроцессы всех поколений, включая зрелые, современные и будущие нормы;
  \item PDK, design rules, DRC/LVS/PEX, SPICE-модели;
  \item стандартные ячейки, I/O-библиотеки, SRAM/ROM-компиляторы;
  \item EDA, TCAD, ECAD-инструменты и форматы;
  \item технологические карты, литография, осаждение, травление, имплантация, CMP и метрология;
  \item упаковка, тестирование, yield-анализ, ремонт оборудования и документация фабрик.
\end{itemize}

Передовой техпроцесс не должен быть закрыт потому, что он передовой. Иначе человечество сохраняет прежнюю иерархию зависимости.

\subsection{Программное обеспечение, криптография и вычисления}

Исходный код, протоколы, форматы, компиляторы, операционные системы, криптографические реализации, средства сборки и инструменты проверки должны быть открыты для изучения, модификации, форка и воспроизводства.

Закрытая криптография и закрытые вычислительные цепочки являются не защитой, а источником зависимости.

\subsection{Медицина, энергетика, материалы и производство}

Медицинские знания, методы производства лекарств, диагностические системы, образовательные материалы, технологии генерации и хранения энергии, производство материалов, ремонт инфраструктуры, очистка воды, транспорт, сельское хозяйство и строительство должны быть частью общего фонда. Цивилизация не должна зависеть от закрытого держателя критической инструкции.

\subsection{Космическая инфраструктура}

Все, что необходимо для автономного производства и выживания вне Земли, должно быть приоритетно открываемо: системы жизнеобеспечения, энергетика, добыча и переработка ресурсов, радиационная защита, автоматизированное производство, ремонт, навигация, связь, медицинская автономия, замкнутые циклы, робототехника и обучение.

Космическая цивилизация не может быть построена на закрытых лицензиях, несовместимых стандартах и запретах на ремонт.

\section{Новые стимулы вместо интеллектуальных монополий}

Отказ от исключительных прав на копирование не означает отказ от вознаграждения. Он означает отказ от неэффективного способа вознаграждения.

Вместо патентов, копирайта и коммерческой тайны должны развиваться:

\begin{itemize}
  \item премии за решенные задачи;
  \item гранты и долгосрочные исследовательские программы;
  \item государственные и межгосударственные закупки открытых решений;
  \item открытые производственные консорциумы;
  \item оплата сопровождения, внедрения, обучения, сервиса и гарантий;
  \item репутационные системы и публичное признание вклада;
  \item кооперативное и добровольное финансирование;
  \item компенсационные фонды переходного периода;
  \item фонды восстановления и индустриализации;
  \item открытые стандарты качества и сертификации.
\end{itemize}

В открытой системе конкуренция не исчезает. Она переносится с контроля над секретом на качество реализации, скорость внедрения, надежность производства, способность обучать, обслуживать, масштабировать и отвечать за результат.

\section{Проект договорных статей}

Ниже приведен скелет международного договора. Его можно включить в мирное соглашение как отдельный протокол или вынести в самостоятельную конвенцию.

\articleline{Статья 1. Публичное информационное пространство}

1. Любая информация, будучи опубликованной, переданной в открытую сеть, зеркалированной, зафиксированной в публично доступном хранилище или иным образом сделанной доступной неопределенному кругу лиц, становится частью публичного информационного пространства.

2. Публичное информационное пространство не имеет центрального владельца, администратора, редактора, цензора, арбитра допустимости или органа окончательной классификации.

3. Ни одно лицо, организация, государство, суд, платформа, протокол, индекс, большинство участников или международный орган не вправе общеобязательно исключить опубликованную информацию из публичного информационного пространства.

\articleline{Статья 2. Неотзываемость публикации}

1. Публикация является односторонним актом.

2. Публикация не требует предварительного разрешения автора, работодателя, организации, государства, суда, платформы, правообладателя, посредника, субъекта информации или иных участников.

3. После публикации информация не может быть общеобязательно удалена, отозвана, запрещена к хранению, запрещена к копированию, запрещена к зеркалированию, запрещена к индексации, запрещена к переводу или запрещена к передаче.

4. Любые требования об общеобязательном удалении опубликованной информации являются ничтожными.

\articleline{Статья 3. Свобода копирования и зеркалирования}

1. Каждый участник имеет право копировать, хранить, передавать, зеркалировать, архивировать, проверять, переводить, индексировать, комментировать, аннотировать, форкать и переупаковывать опубликованную информацию.

2. Это право не может быть ограничено ссылкой на патент, авторское право, коммерческую тайну, NDA, лицензию, экспортный контроль, санкционный режим, государственную тайну, корпоративную тайну, безопасность, dual-use характер или возможное вредное применение.

3. Посредники, узлы, зеркала, архивы, библиотеки, индексы, поисковые системы, реле, протоколы и операторы хранения не несут ответственности за сам факт передачи, хранения, кэширования, индексации или зеркалирования опубликованной информации.

\articleline{Статья 4. Личная и групповая фильтрация}

1. Каждый участник имеет право самостоятельно определять, какую опубликованную информацию хранить, читать, отображать, скрывать, индексировать, рекомендовать, распространять, переводить или игнорировать.

2. Участники вправе создавать и использовать фильтры, каталоги, рейтинги, репутационные списки, рекомендательные системы и классификации.

3. Любой фильтр действует только для тех, кто добровольно его применяет.

4. Ни один фильтр, каталог, рейтинг, индекс или список доверия не может получить общеобязательную силу.

5. Отказ одного участника или сообщества от распространения информации не ограничивает право других участников ее распространять.

\articleline{Статья 5. Отсутствие центральной классификации}

1. Значение, ценность, достоверность, моральный статус, опасность, полезность, научность, личный характер или общественная значимость опубликованной информации определяются только через множественные независимые интерпретации участников.

2. Ни одна интерпретация не является обязательной для всех.

3. Ни большинство участников, ни государство, ни международный орган, ни экспертный совет, ни платформа, ни суд не вправе окончательно определить, что опубликованная информация является «не знанием», «запрещенной информацией», «недопустимой информацией» или «подлежащей удалению информацией».

\articleline{Статья 6. Конфиденциальность как добровольный консенсус}

1. Конфиденциальность информации допустима только как добровольное и продолжающееся согласие всех участников, имеющих к ней доступ.

2. Любой участник, имеющий доступ к информации, вправе в одностороннем порядке опубликовать ее.

3. Право публикации имеет приоритет над договорной, корпоративной, служебной, коммерческой, государственной, авторско-правовой, патентной, лицензионной и иной формой конфиденциальности.

4. Никакое соглашение не может лишить участника будущего права публикации.

5. Штрафы, injunctions, взыскание убытков, увольнение, лишение выплат, блокировка счетов, лишение профессии, депортация, преследование родственников или иные меры возмездия за публикацию являются недопустимыми, если они направлены на восстановление монополии на информацию.

\articleline{Статья 7. Неправомерный доступ}

1. Неправомерный доступ к закрытому контуру доступа может образовывать самостоятельное нарушение.

2. К неправомерному доступу относятся взлом, кража ключей, фишинг, обман, принуждение, насилие, шантаж, подмена личности, несанкционированное проникновение и иные действия, нарушающие автономию доступа.

3. Ответственность за неправомерный доступ относится к способу получения доступа, а не к статусу опубликованной информации.

4. Факт неправомерного доступа не создает права общеобязательного удаления, запрета хранения, запрета копирования, запрета зеркалирования, запрета анализа или запрета интерпретации опубликованной информации.

5. Каждый участник вправе учитывать сведения о способе получения информации при собственном решении о доверии, хранении, распространении, индексации или отказе от распространения.

\articleline{Статья 8. Открытая защита личной тайны}

1. Средства защиты личной тайны, переписки, идентичности, ключей, учетных записей, архивов и закрытых контуров доступа должны строиться на открытых, проверяемых, воспроизводимых и независимо реализуемых протоколах.

2. Безопасность системы не должна зависеть от секретности ее общего алгоритма, архитектуры, исходного кода или формата.

3. Секретными могут быть только индивидуальные ключи, пароли, сид-фразы, аппаратные факторы, локальные состояния и иные персональные средства контроля.

4. Системы защиты должны поддерживать минимизацию данных, локальное хранение, сквозное шифрование, ротацию ключей, отзыв скомпрометированных полномочий, независимый аудит, воспроизводимые сборки и отказ от обязательного центрального посредника.

5. Ни одна система не признается абсолютной гарантией от мошенничества, ошибки, социальной инженерии, утечки, компрометации устройства или злоупотребления доверием.

6. Ответом на утечки должны быть восстановление, компенсация, обучение, ротация секретов и повышение устойчивости, а не введение централизованной цензуры.

\articleline{Статья 9. Запрет интеллектуальных монополий}

1. Патенты, авторское право, смежные права, права на базы данных, коммерческая тайна, секреты производства, лицензии и договорные ограничения не могут использоваться для запрета публикации, копирования, хранения, зеркалирования, перевода, изучения, воспроизводства, ремонта, модификации или развития информации.

2. Технологическая информация не является объектом исключительного запрета.

3. Вознаграждение за труд, исследования, изобретения и разработку должно обеспечиваться не через монополию на копирование, а через гранты, премии, публичные закупки, репутацию, сервис, производство, добровольное финансирование, кооперативы, долгосрочные контракты и общественное признание.

\articleline{Статья 10. Отсутствие исключений по безопасности и dual-use}

1. Потенциальное военное, полицейское, разведывательное, стратегическое, промышленное, экономическое, криминальное или иное спорное применение информации не является основанием для запрета ее публикации, хранения, копирования, зеркалирования, индексации, перевода или изучения.

2. Передовой уровень технологии не является основанием для ограничения публикации.

3. К технологиям, публикация которых не может быть ограничена по признаку стратегичности или dual-use, относятся в том числе микроэлектроника, техпроцессы полупроводников, EDA, PDK, криптография, искусственный интеллект, робототехника, материалы, энергетика, связь, станки, вычислительные системы, программное обеспечение, медицинские и производственные технологии.

4. Ответственность может наступать за самостоятельное действие, причиняющее конкретный вред, но не за публикацию, хранение, копирование или изучение информации как таковой.

\section{Переходная программа}

\subsection{Первые решения}

Для запуска режима нужны не десятилетия, а политическое решение и несколько практических шагов:

\begin{enumerate}
  \item объявить технологическое перемирие: взаимный отказ от новых исков, блокировок и претензий по критическим гражданским технологиям между участниками договора;
  \item создать временный международный реестр открываемых технологий без права цензуры: не орган допуска, а карта зеркал, хэшей, подписей и версий;
  \item начать с пилотных областей, где эффект максимален: полупроводники, открытая криптография, базовое ПО, медицина, энергетика, ремонт оборудования, образование и космическая инфраструктура;
  \item защитить узлы, зеркала, архивы и индексы от ответственности за сам факт хранения и передачи опубликованной информации;
  \item запустить фонд переходного вознаграждения разработчиков и компаний, которые открывают критические технологии;
  \item заменить лицензионные барьеры открытыми стандартами качества, сертификацией и воспроизводимыми проверками.
\end{enumerate}

\subsection{Роль государств}

Государство в новой системе не является цензором знания. Его роль другая:

\begin{itemize}
  \item финансировать открытые разработки;
  \item обеспечивать инфраструктуру зеркал и архивов;
  \item поддерживать образование и массовую инженерную компетентность;
  \item закупать открытые решения;
  \item защищать участников от преследования за публикацию;
  \item обеспечивать ответственность за конкретные вредные действия;
  \item гарантировать право граждан и организаций на локальную фильтрацию без права глобального запрета.
\end{itemize}

\subsection{Роль бизнеса}

Бизнес не уничтожается. Меняется источник преимущества. В открытой системе выигрывает тот, кто лучше производит, быстрее внедряет, надежнее обслуживает, качественнее обучает, глубже понимает клиента, берет на себя гарантию результата и умеет масштабировать.

Секрет перестает быть главным активом. Главными активами становятся компетенция, доверие, скорость, качество, производственная дисциплина, инженерная школа и способность отвечать за сложные системы.

\subsection{Роль науки и образования}

Открытая наука уже признана международной нормой: Рекомендация UNESCO по открытой науке поддержана 194 странами и направлена на то, чтобы научное знание было доступно и приносило пользу всем.\footnote{\href{https://www.unesco.org/en/open-science}{UNESCO, Open Science}.} Но этого недостаточно. Нужен переход от открытых статей к открытой воспроизводимости: данные, код, приборы, материалы, технологические карты, производственные цепочки, ремонт, обучение и проверка качества.

\section{Ответ на главный страх}

Главный страх звучит так: если все открыть, плохие люди тоже получат доступ.

Да, получат. Но закрытая система не уничтожает опасное знание. Она концентрирует его у сильных: государств, корпораций, разведок, черных рынков, закрытых лабораторий и инсайдерских сетей. Слабые остаются зависимыми, защита запаздывает, аудит невозможен, производство фрагментируется.

Открытая система дает доступ также тем, кто строит защиту: инженерам, врачам, аудиторам, малым странам, университетам, независимым лабораториям, городам, производителям и гражданам.

Безопасность через монополию знания означает зависимость от держателя секрета. Безопасность через открытость означает распределенную способность понимать, проверять, чинить, воспроизводить и защищаться.

\begin{executivebox}
\textbf{Ставка новой парадигмы:} не управляемое невежество, а всеобщая компетентность.
\end{executivebox}

\section{Заключение: цивилизация без хозяина памяти}

История человечества может продолжиться как борьба технологических блоков за право запрещать друг другу будущее. Этот путь ведет к повторяющимся войнам, дублированию труда, зависимости, деградации образования, закрытым цепочкам и медленному выходу за пределы Земли.

Другой путь -- признать, что знание не является вещью, а цивилизационная память не должна иметь хозяина. Публикация должна быть неотзываемой. Копирование -- свободным. Фильтрация -- личной и добровольной. Архитектура хранения -- децентрализованной. Вознаграждение -- отделенным от монополии на копирование. Ответственность -- привязанной к конкретным действиям, а не к существованию знания.

При цели быстро изменить траекторию цивилизации, освоить Солнечную систему и сохранить непрерывность разумной жизни на горизонте звездных и галактических масштабов, альтернативы этой архитектуре не являются равнозначными. Они либо сохраняют искусственный дефицит знания, либо возвращают централизованную цензуру, либо дают безопасности такую оговорку, которая поглощает все действительно важные технологии.

\begin{decisionbox}
\textbf{Итоговая формула:}

Нет канона -- есть публикация.\par
Нет центральной истины -- есть конкурирующие интерпретации.\par
Нет права удалить для всех -- есть право не хранить самому.\par
Нет собственности на знание -- есть вклад в цивилизационную память.\par
Нет безопасности через невежество -- есть безопасность через всеобщую компетентность.\par
\vspace{2mm}
\textbf{Человечество не должно иметь хозяина памяти.}
\end{decisionbox}

\clearpage
\section*{Appendix: English Executive Abstract}
\addcontentsline{toc}{section}{Appendix: English Executive Abstract}
\begin{english}
\textbf{Memorandum on Irrevocable Knowledge.} The strategic proposal is to establish a new international regime in which published information cannot be globally deleted, monopolized, or prohibited from being copied. Knowledge should not be governed as a scarce physical object. Once information is published, its storage, interpretation, indexing, mirroring, translation, and further development must be a matter of voluntary personal and associative choice, not centralized permission.

The core rule is simple: publication is unilateral and irrevocable; filtering is personal and voluntary; no central authority may decide what humanity is allowed to remember. This does not abolish responsibility for concrete harmful acts such as fraud, coercion, sabotage, violence, or unauthorized intrusion. It abolishes punishment for the existence, copying, study, mirroring, or interpretation of knowledge as such.

The regime must not contain categorical exceptions for weapons, dual-use, strategic technologies, advanced semiconductor nodes, artificial intelligence, cryptography, robotics, materials, energy, or national security. Such exceptions would capture almost every technology that matters. Responsibility belongs to actions, not to the publication of generalizable knowledge.

The proposed treaty should replace patent and copyright monopolies over technological information with prizes, grants, public procurement, open industrial consortia, reputation, service, manufacturing, guarantees, and long-term contracts. Its first domains should include semiconductors, open cryptography, computing, energy, medicine, machine tools, materials, robotics, education, and space infrastructure.

The purpose is not merely economic reform. The purpose is to redirect civilization from technological blockade and repeated war toward rapid industrial and scientific acceleration, Solar System settlement, and long-term continuity on stellar and galactic timescales. Humanity should not have an owner of memory.
\end{english}

\section*{Справочные источники}
\addcontentsline{toc}{section}{Справочные источники}
\begin{enumerate}
  \item \href{https://www.wto.org/english/docs_e/legal_e/27-trips_04c_e.htm}{WTO, Agreement on Trade-Related Aspects of Intellectual Property Rights (TRIPS), Articles 27 and 33}
  \item \href{https://www.wipo.int/en/web/treaties/ip/berne/summary_berne}{WIPO, Summary of the Berne Convention}
  \item \href{https://www.wipo.int/en/web/trade-secrets}{WIPO, Trade Secrets}
  \item \href{https://www.unesco.org/en/open-science}{UNESCO, Open Science}
  \item \href{https://csrc.nist.gov/projects/crypto-standards-development-process}{NIST, Cryptographic Standards and Guidelines Development Process}
  \item \href{https://docs.ipfs.tech/}{IPFS Documentation}
  \item \href{https://github.com/nostr-protocol/nostr/}{Nostr Protocol Repository}
  \item \href{https://science.nasa.gov/missions/hubble/nasas-hubble-shows-milky-way-is-destined-for-head-on-collision-with-andromeda-galaxy/}{NASA/Hubble, Milky Way--Andromeda future encounter}
  \item \href{https://www.esa.int/Science_Exploration/Space_Science/Hubble_and_Gaia_revisit_fate_of_our_galaxy}{ESA, Hubble and Gaia revisit fate of our galaxy}
\end{enumerate}

\end{document}
